\section{Linear Programming}

	We are interested in solving constrained linear optimization problems that can be described in the \hyperref[opt-def:primal_dual]{primal/dual formulation}. The theory is all established in my optimization notes, so we will apply them to various problems such as: 
  \begin{enumerate}
    \item Shortest path: select a set of edges of minimum cost such that $s$ and $t$ are connected. 
    \item Minimum Spanning Tree: select a set of edges of minimum cost such that all vertices are connected. 
    \item Maximum Flow: select a flow on every edge such that capacity and balance constraints are respected. 
  \end{enumerate} 
	These are all instances of \textit{linear programming}. 

  \begin{theorem}[Optimality of LP]
    An optimal solution to a linear program can be obtained in weakly polynomial time in the number of variables, constraints, and poly-logarithmic in the constants in the LP.\footnote{Obtaining a strongly polynomial algorithm for LPs remains a major open problem.} 
  \end{theorem} 

  This is very similar to integer programs, which solve for integer solutions rather than real, but in general we don't know of algorithms that can solve integer programs in polynomial time. However, integer programs can be ``relaxed'' to linear programs by removing the restriction on variables being integers. However, the optimal solution for the relaxed LP may be different from that of the original LP. 

  \begin{example}[Max Flow as LP]
    
  \end{example}

  \begin{example}[Bipartite Matching as LP]
    
  \end{example} 

  \begin{example}[Vertex Cover]
    You are given an undirected graph $G = (V, E)$. A vertex cover is a subset of vertices $S \subseteq V$ such that every edge in $E$ has at least one endpoint in $S$. That is, for each edge $(u, v) \in E$, either $u \in S$ or $v \in S$ (or both). The size of a vertex cover $S$ is the number of vertices in $S$, $|S|$. The minimum vertex cover is the vertex cover with the minimum size.
    \begin{enumerate}
      \item Write an integer program to compute the minimum vertex cover in $G$.  Clearly define all variables and provide a brief English explanation of the objective and each constraint.
      \item Relax your program from part $a$ to a linear program, and take it's dual. Interpret your result as another algorithmic optimization problem.
    \end{enumerate} 

    The relaxed primal LP problem is finding 
    \begin{equation}
      \min \sum_{v \in V} x_v \text{ subject to } \begin{cases} 
        x_u + x_v \geq 1 & \text{ for all } (u, v) \in E \\
        x_u \geq 0 & \text{ for all } u \in V
      \end{cases}
    \end{equation}
    To introduce the dual problem, we can introduce the variables for each edge $x_{uv} \geq 0$. Therefore, we can take each constraint $x_u + x_v \geq 1$ and multiply it by $x_{uv}$ to get $x_{uv} (x_u + x_v) \geq x_{uv}$ for all $(u, v) \in E$. Therefore, summing it up, we get 
    \begin{equation}
      \sum_{(u, v) \in E} x_{uv} (x_u + x_v) \geq \sum_{(u, v) \in E} x_{uv}
    \end{equation}
    Now we rephrase the left summation. Summing over all edges is the same as summing over all nodes, and for each node summing over all adjacent edges to that node, and so we can get 
    \begin{equation}
      \sum_{v \in V} \bigg( \sum_{(u, v) \in E} x_{uv} \bigg) x_v \geq \sum_{(u, v) \in E} x_{uv} 
    \end{equation}
    To make the LHS a lower bound of $\sum_{v \in V} x_v$, we can introduce the constraints that all coefficients are $\leq 1$, so we have the dual constraints 
    \begin{equation}
      \sum_{(u, v)} x_{uv} \leq 1 \text{ for all } v \in V
    \end{equation}
    Since $\sum_{(u, v) \in E} x_{uv}$ is a lower bound, we want to maximize the lower bound and so our dual optimization problem is 
    \begin{equation}
      \max \sum_{(u, v) \in E} x_{uv}
    \end{equation}
  \end{example}
